\documentclass[11pt]{article}
\usepackage{arxiv-preprint}

\title{Filler-Augmented Training Extends Depth\\ and Prevents State Divergence\\ in Fast-Weight Recurrent Sequence Models}

\author[1]{Voidwolf}
\affil[1]{Independent \\ \texttt{research@voidwolf.space}}

\date{August 2026}

\begin{document}
\maketitle

\begin{abstract}
Fast-weight recurrent sequence models with content-addressed memory
can compose multi-hop associative retrievals inside a single forward
pass, but training at moderate depth reveals three coupled failures:
\emph{query-conditioned formation}, in which a fraction of seeds
converge to a low-rank readout that collapses when any filler
separates the last fact from the query; \emph{state divergence}, in
which the residual-stream norm grows exponentially past a
sequence-length threshold and the readout collapses onto whichever
token direction the explosion aligns with; and a
\emph{baseline depth wall} beyond which no healthy seed forms at all.
We show that a single training-time modification addresses all
three failures simultaneously: insert a uniform-random count of
filler tokens between the last fact and the query, with a
warm-up-and-ramp schedule and a two-task joint loss. Across trained hop
counts $K_{\mathrm{hops}} \in \{3, 4, 5, 6\}$, the recipe drops the
mean query-conditioning diagnostic from above the state-persistent
boundary to below it, eliminates seeds whose residual-stream norm
exceeds $10^3$ at $N=2000$ filler positions, and restores a $4/5$
pass rate at the $K_{\mathrm{hops}}=6$ cell where the baseline recipe
produced $0/3$ healthy seeds. The intervention is scaffold-agnostic:
an explicitly two-layer variant of the same architecture solves the
$K_{\mathrm{hops}}=3$ task under the same recipe with target
accuracy $0.992$ across three seeds. We report a known nuance: at depth the recipe biases the model
toward the filler-present distribution, so pass-seeds increasingly
need filler at evaluation to reach their headline accuracy. We
close by placing the recipe inside the broader landscape of
training-time interventions for recurrent recall.
\end{abstract}

% ======================================================================
\section{Introduction}
\label{sec:intro}

Fast-weight recurrent sequence models maintain an outer-product
associative memory as their hidden state, updating it in place per
token. The family runs from \citet{schlag2021linear}, which cast
linear transformers as fast-weight programmers, through the modern
DeltaNet lineage \citep{yang2024deltanet,yang2024gated} and its
self-referential branch \citep{irie2022srwm}. On single-hop associative recall, the family closes much of
the recall-capacity gap between Mamba-style state-space models
\citep{gu2024mamba,dao2024mamba2} and softmax attention
\citep{arora2023zoology,arora2024based}. On \emph{multi-hop}
composition, where the answer to a query depends on chaining several
stored $(k, v)$ pairs, fast-weight models matter for a further
reason: the fast-weight state, unlike Mamba's diagonal SSM state,
admits a write that composes state-dependent addresses inside a
single forward pass. A recent extension of the DeltaNet
lineage supplies exactly such a write and composes hop-chains of
length $K_{\mathrm{hops}} \le 5$ inside a fixed three-layer stack at
the trained per-hop bank capacity.

At $K_{\mathrm{hops}} = 6$, the baseline training recipe fails
categorically: across three replicated seeds, one lands near chance,
one appears to converge but is measured to hold the answer in a
representation that collapses under a single filler token, and one
diverges numerically. The failure is not "the trained network is not
accurate enough"; it is three distinct pathologies, all present, all
mixed within the seed distribution. We call them
\emph{query-conditioned formation}, \emph{state divergence}, and the
\emph{depth wall}. They admit separate mechanistic descriptions
(Sec.~\ref{sec:failures}) but they respond, jointly, to one training
recipe: inserting a uniform-random number of filler tokens between the
last stored fact and the query, ramped in over a warm-up window.

\paragraph{Contributions.} This paper isolates and reports that
recipe. Specifically we contribute:

\begin{enumerate}
\item A formal description of the filler-augmented training procedure
      (schedule, loss structure, and hyperparameter defaults) that is
      architecture-generic within the fast-weight family
      (Sec.~\ref{sec:method}).
\item Evidence that the recipe restores a healthy pass rate at the
      $K_{\mathrm{hops}}=6$ depth cell where the baseline recipe
      produces zero healthy seeds (Sec.~\ref{sec:depth-wall}).
\item Evidence that the same recipe, at the same hyperparameters,
      drops the query-conditioning diagnostic below the
      state-persistent boundary across
      $K_{\mathrm{hops}} \in \{3, 4, 5, 6\}$ (Sec.~\ref{sec:formation})
      and eliminates state divergence across the same range
      (Sec.~\ref{sec:divergence}).
\item Two ablations: (a) the recipe is scaffold-agnostic, in that an
      explicitly two-layer variant of the architecture solves under
      the same procedure with mean target accuracy $0.992$ across
      three seeds (Sec.~\ref{sec:ablations}); (b) an occlusion-invariance
      auxiliary loss steers formation into the same mode, providing
      a second training-time route to the same effect
      (Sec.~\ref{sec:ablations}).
\item A characterisation of the recipe's principal cost: at
      $K_{\mathrm{hops}} \ge 4$ the recipe biases the model toward
      the filler-present distribution, so pass-seeds solve strictly
      better with filler than without (Sec.~\ref{sec:filler-dep}).
\end{enumerate}

\paragraph{Position.} The recipe belongs to the catalogue of
training-time interventions that shape which representation
gradient descent commits to when several representations solve the
training distribution equally well
\citep{wang2024grokked,olsson2022inductionheads,liu2023towards}. The
distinctive property here is that a single intervention resolves
three empirically separable failure modes at once, at the cost of
one hyperparameter, and does so through a mechanism that is
mechanistically legible: the intervention forces the readout to
operate on the fast-weight state rather than the query's local
receptive field, and it forces the write on filler tokens to be
near-zero. Both of those consequences fall out of the SEP-write
geometry the recipe induces.

% ======================================================================
\section{Background}
\label{sec:background}

We assume familiarity with fast-weight recurrent layers; readers
who want a survey are referred to
\citet{schlag2021linear,yang2024deltanet,irie2022srwm}. This section
fixes notation, states the retrieval task the paper's experiments
run on, and names the three failure modes the recipe is designed to
address.

\subsection{The mixer}

At each position $t$ the layer computes query, key, value, and
gate projections
$q_t = W_q x_t$, $k_t = W_k x_t$, $v_t = W_v x_t$,
$\beta_t = \sigma(W_\beta x_t)$, then a state-dependent effective
key $k^{\mathrm{eff}}_t = k_t + M_{t-1} k_t$, updates
\begin{equation}
  M_t \;=\; M_{t-1} \;+\; \beta_t \cdot k^{\mathrm{eff}}_t \, v_t^\top,
  \label{eq:write}
\end{equation}
and reads $y_t = M_t^\top q_t$. The state
$M_t \in \mathbb{R}^{d_k \times d_v}$
is an outer-product associative memory; the effective-key term
$M_{t-1} k_t$ is a state \emph{lookup} of the raw key, retrieving
the stored key whose associated value most resembles the current
input's key, so the write address is a linear function of what is
already stored under $k_t$ at the moment of the write.\footnote{This is the minimal
address-side residual in the sense discussed in
\citet{irie2022srwm,irie2021going}. It keeps the recurrence affine in
$M$, which is what allows the layer to be chunk-parallelised at the
DeltaNet complexity class \citep{yang2024deltanet}.} The layer is
placed inside a standard Mamba-style scaffold
\citep{gu2024mamba,dao2024mamba2}: input projection, grouped
convolution, SiLU, $z$-branch gate, output projection. Throughout the
paper we use $L=3$ layers, $d_{\mathrm{model}} = 256$,
$d_k = d_v = 32$, and bf16 autocast unless otherwise stated.

\subsection{The retrieval task}

The task is $K_{\mathrm{hops}}$-hop factorised associative retrieval,
a multi-hop generalisation of the single-hop associative-recall
tasks in \citet{arora2023zoology,arora2024based} developed for the
companion mechanism paper \citep{fenrir2026}. Each row is a concatenation of
$K_{\mathrm{hops}} - 1$ intermediate fact banks (each holding $K_i$
$(k, v)$ pairs) and a final bank (holding $K_N$ pairs), separated by
SEP tokens, followed by a query token and a target token to be
predicted. Bridge assignments between successive banks are drawn as
random permutations, so a "confused-entity" traversal cannot land on
the correct target by construction. Chance is $1/K_N$; we hold
$K_i = 4$ for $i < N$ and $K_N = 8$, giving chance $0.125$. This is
the "$K_1=4/K_N=8$ matched-difficulty" setting inherited from
earlier work on the mechanism; it is the shape at which a positive
result cannot be reached by simply outputting some value from a
small set.

\subsection{Three failure modes at depth}
\label{sec:failures}

Under the baseline training recipe (no filler, joint half-half loss
on a $K_1 = 1$ probe and the $K_{\mathrm{hops}}$ target,
$15{,}000$ steps, batch $64$), the trained model fails at depth in
three distinct ways.

\paragraph{(F1) Query-conditioned formation.} A fraction of seeds
(one third to one half in our replications) converge to a solution
whose readout depends on local content adjacent to the query token
rather than on the fast-weight state. Inserting even a handful of
filler tokens between the last fact and the query collapses the
prediction. The failure has a clean training-time diagnostic:
QTOK-flip \citep{fenrir2026}, the fraction of the prediction
distribution that changes when the query token is half-occluded,
above $0.15$ at step $5000$ predicts the mode at step $15000$.
Section~\ref{sec:formation} treats this diagnostic in the results.

\paragraph{(F2) State divergence.} A separate fraction of seeds
(growing with $K_{\mathrm{hops}}$) produce a readout that
appears healthy on the training-length distribution but whose
residual-stream norm grows exponentially at inference length. On
audited baseline seeds, the norm at $N = 500$ inserted filler
positions can already reach $10^5$; at $N = 2000$ it can exceed
$10^{12}$ (measured on rig-056 $K_{\mathrm{hops}}=5$ seeds via a
post-hoc divergence audit). The readout collapses to whichever vocabulary token
happens to align with the direction of the explosion. Different
divergent seeds diverge in different directions, so the collapse is
seed-specific rather than model-specific.

\paragraph{(F3) The depth wall.} Under the baseline recipe, hop
counts $K_{\mathrm{hops}} \in \{2, 3, 4, 5\}$ are solved with mean
accuracy above $0.94$ across three-seed replications. At
$K_{\mathrm{hops}} = 6$, three replicated seeds all fail: one lands
at chance with mild query-conditioning, one lands at $0.86$ target
but with QTOK-flip $= 0.77$ and NaN state on filler probes (a
degenerate solve through the query window rather than the state),
and one lands at chance with a state that diverges by $10^9$-scale
norm at $N = 2000$. The failure is not a matter of insufficient
budget; each of the three failures reflects a different local
optimum the baseline recipe steers seeds into as the mechanism has
to compose one more hop than the state's write geometry easily
supports.

% ======================================================================
\section{The recipe}
\label{sec:method}

The intervention has three components: an insertion, a schedule, and
a joint loss.

\subsection{Insertion}

For each training row containing $K_{\mathrm{hops}}$-hop target
structure, before appending the query token, we insert $n$ copies of
the SEP token between the last bank's final fact and the query,
where
\[
  n \;\sim\; \mathrm{Uniform}\{0, 1, \dots, n_{\max}(t)\}.
\]
One $n$ is drawn per batch (a single filler length applies to all
rows within a batch), redrawn every step. $n_{\max}(t)$ is a
function of the training step $t$, described next.

\subsection{Schedule}

$n_{\max}(t)$ starts at $0$, holds at $0$ across a warm-up window
$[0, T_{\mathrm{warm}}]$, then ramps linearly to a target cap
$n^\star$ across $[T_{\mathrm{warm}}, T_{\mathrm{warm}} + T_{\mathrm{ramp}}]$,
and holds at $n^\star$ thereafter:
\[
  n_{\max}(t) \;=\; n^\star \cdot \mathrm{clip}\!\left(
    \frac{t - T_{\mathrm{warm}}}{T_{\mathrm{ramp}}},\; 0,\; 1
  \right).
\]
The warm-up matters: the mechanism needs to form on the
short-sequence distribution before it can be pulled onto the
filler-present distribution. Skipping warm-up produces seeds that
never form the mechanism at all. Throughout the paper we use
$T_{\mathrm{warm}} = 2000$ steps, $T_{\mathrm{ramp}} = 2000$ steps,
and $n^\star = 200$. These defaults have not been tuned per-depth;
Sec.~\ref{sec:ablations} shows they carry the $K_{\mathrm{hops}}=6$
cell without adjustment.

\subsection{Loss}

Every step draws two batches of equal size: a $K_1 = 1$ probe batch
with no filler, and a $K_{\mathrm{hops}}$ target batch with filler
drawn as above. The total loss is the equal-weight average of the
two per-token cross-entropies. The probe batch keeps the write
geometry on short sequences aligned with the target task's write
geometry; ablating it (target-only loss with filler) produces seeds
that solve the filler-present distribution but not the filler-free
distribution and vice versa.

\subsection{What it changes about the write}

The recipe's effect can be read out of Eq.~\ref{eq:write}: on filler
(SEP) positions the gate $\beta_t$ must not write meaningful content
into $M$, because the target loss punishes any state accumulation
that survives up to the query position. This has two consequences.
First, the readout $M_T q_T$ at the query position is forced to
retrieve from what was stored at the fact positions, not from
whatever the last few positions before the query happened to
contain. That fixes the query-conditioned failure (F1) by
construction.
Second, the state stops accumulating on filler tokens, which fixes
the state divergence failure (F2) at its geometric cause: the
divergence traced by \citet{fenrir2026} is precisely a runaway
accumulation of non-negligible writes on many SEP positions.

% ======================================================================
\section{Experimental setup}
\label{sec:setup}

We report five-seed replications for each $K_{\mathrm{hops}}$ cell in
$\{3, 4, 5, 6\}$, plus a three-seed replication for the two-layer
ablation. Training runs $15{,}000$ steps at batch $64$ with the
AdamW optimiser (weight decay $0.01$, $\beta = (0.9, 0.95)$),
cosine LR schedule with base LR $3 \cdot 10^{-4}$, $200$-step
warmup, and floor $3 \cdot 10^{-5}$ ($10\%$ of base). Gradients
are clipped to unit norm before each step. Autocast is bf16. All
experiments run on an AMD Radeon RX 7900 XTX (RDNA3, 24\,GB, via
ROCm 7.2); a chunked implementation of the mixer
\citep{yang2024deltanet,fenrir2026} produces a $4.2$x wallclock
speedup on the $K_{\mathrm{hops}}=2$ shape and $6.0$x on
$K_{\mathrm{hops}}=3$; loop-parity holds at $K_{\mathrm{hops}}=5$
to $1.5 \cdot 10^{-6}$ relative error.

\paragraph{Metrics.}

\begin{itemize}
\item \textbf{Target accuracy at N=200 fillers} (target$_{N=200}$):
      the primary headline. Evaluated on $2000$ held-out rows
      constructed identically to training but with fixed $n = 200$
      filler tokens inserted.
\item \textbf{Target accuracy at N=0 fillers} (target$_{N=0}$):
      how the seed behaves on the filler-free distribution.
      Difference from target$_{N=200}$ diagnoses
      filler-dependence (Sec.~\ref{sec:filler-dep}).
\item \textbf{QTOK-flip}: fraction of predictions that change when
      the query token is half-occluded, at a fixed filler length
      $N = 200$. Above $0.15$ = query-conditioned; at or below
      $0.15$ = state-persistent \citep{fenrir2026}. Passing seeds
      under the recipe cluster well below the boundary (typically
      $\le 0.03$).
\item \textbf{div\_norm@N=2000}: residual-stream norm at position
      $\mathrm{seq\_len} - 1$ when evaluated at $N = 2000$ inserted
      fillers. Above $10^3$ = diverging; below $10^2$ = healthy
      \citep{fenrir2026}.
\end{itemize}

\paragraph{Verdict rules.} A seed \emph{passes} its cell if
target$_{N=200} \ge 0.80$, QTOK-flip $\le 0.15$
(state-persistent), and div\_norm@N=2000 $< 10^3$. A cell passes
if $\ge 4/5$ seeds pass; a cell is a strong pass if all seeds
pass. The rules are pre-registered from the earlier design
document; the $0.80$ threshold sits well above chance $0.125$.

% ======================================================================
\section{Results}
\label{sec:results}

The main empirical claim: one recipe, unchanged across four hop
counts, restores target accuracy, formation mode, and numerical
stability on all four cells. The same recipe, at the same
hyperparameters, breaks the $K_{\mathrm{hops}} = 6$ baseline wall.

\subsection{Depth wall broken at K\_hops=6}
\label{sec:depth-wall}

Table~\ref{tab:depth-wall} places the recipe's $K_{\mathrm{hops}}=6$
result next to the baseline's. The baseline recipe produces zero
healthy seeds across three replications: one at chance with mild
query-conditioning, one nominally solving the target but with a
readout that is entirely query-adjacent (QTOK-flip $= 0.77$), and
one with a residual-stream norm of $1.19 \cdot 10^9$ at $N = 2000$.
The filler recipe produces four passing seeds out of five with mean
target$_{N=200} = 0.867$ across the passers, QTOK-flip $\le 0.02$
for all five, and div\_norm$_{N=2000} \le 78$ on the passers. The
one failing seed under the recipe (s3) sits near chance
(target$_{N=200} = 0.312$) with an elevated div\_norm of $2043$;
formation failure and mild divergence co-occur in the same seed
rather than manifesting as distinct events.

\begin{table}[h]
\centering
\small
\caption{The $K_{\mathrm{hops}} = 6$ cell. Baseline (rig 129) vs
filler recipe (rig 136). $n$ = number of replicated seeds. Passer
mean and QTOK-flip mean computed across seeds meeting the
verdict rule.}
\label{tab:depth-wall}
\begin{tabular}{lcccccc}
\toprule
recipe & $n$ & pass rate & target$_{N=200}$ (mean) & passers mean & QTOK-flip (mean) & div\_norm (mean, passers) \\
\midrule
baseline    & 3 & $0/3$ & $0.450$ (spread $0.625$) & --- & $0.323$ & --- \\
filler      & 5 & $4/5$ & $0.756$                  & $0.867$ & $0.005$ & $76$ \\
\bottomrule
\end{tabular}
\end{table}

\subsection{Formation mode across depths}
\label{sec:formation}

Table~\ref{tab:formation} reports QTOK-flip means across all four
depth cells under the recipe. Every cell sits an order of magnitude
below the $0.15$ query-conditioned threshold; the mean decreases
monotonically with depth. Compared to the baseline, where roughly
$1/3$ of $K_{\mathrm{hops}} = 5$ seeds land query-conditioned by
QTOK-flip \citep{fenrir2026}, the recipe pulls essentially every
seed into the state-persistent basin.

\begin{table}[h]
\centering
\small
\caption{Formation-mode diagnostic (QTOK-flip mean) across depths
under the filler-augmented recipe. Lower is state-persistent; the
boundary is at $0.15$.}
\label{tab:formation}
\begin{tabular}{lcccc}
\toprule
$K_{\mathrm{hops}}$ & 3 & 4 & 5 & 6 \\
\midrule
rig            & 063 & 132 & 133 & 136 \\
seeds          & 5   & 5   & 5   & 5   \\
QTOK-flip mean & $0.033$ & $0.013$ & $0.006$ & $0.004$ \\
\bottomrule
\end{tabular}
\end{table}

\subsection{Divergence suppressed}
\label{sec:divergence}

Under the baseline recipe at $K_{\mathrm{hops}} = 5$, two of three
replicated seeds diverge at $N = 2000$
(div\_norm $\gg 10^5$). Under the recipe, all five replicated
seeds are non-divergent (div\_norm $\le 103$ across the range at
$K_{\mathrm{hops}} = 5$; $\le 78$ across the four passers at
$K_{\mathrm{hops}} = 6$). One seed at $K_{\mathrm{hops}} = 6$
(s3) has an elevated div\_norm of $2043$, but this co-occurs with
that seed's formation failure rather than presenting as a distinct
divergence event.

The reading is that the recipe's SEP-write geometry (forcing
$\beta_t$ near zero on filler positions to satisfy the joint loss)
attacks both the formation-mode failure (F1) and the divergence
failure (F2) at their shared geometric cause: both are
consequences of the write engaging on positions where it should be
idle.

\subsection{Combined result}
\label{sec:combined}

Table~\ref{tab:main} is the paper's headline result. Across
$K_{\mathrm{hops}} \in \{3, 4, 5, 6\}$, the recipe produces mean
target$_{N=200} \ge 0.756$ (or, restricted to passing seeds,
$\ge 0.867$), QTOK-flip mean $\le 0.033$, and eliminates divergence
across $n = 20$ trained seeds with the exception of the one formation
failure at the deepest cell.

\begin{table}[h]
\centering
\small
\caption{Combined recipe results across the trained depth range,
$n = 5$ seeds per cell.}
\label{tab:main}
\begin{tabular}{lccccc}
\toprule
$K_{\mathrm{hops}}$ & target$_{N=200}$ (mean) & QTOK-flip (mean) & passes & non-divergent & filler-dep gap \\
\midrule
3 & $0.995$ & $0.033$ & $5/5$ & $5/5$ & $\sim 0.04$ \\
4 & $0.975$ & $0.013$ & $5/5$ & $5/5$ & $\sim 0.32$ \\
5 & $0.917$ & $0.006$ & $5/5$ & $5/5$ & $\sim 0.50$ \\
6 & $0.756$ (passers $0.867$) & $0.004$ & $4/5$ & $4/5$ & $\sim 0.43$ \\
\bottomrule
\end{tabular}
\end{table}

% ======================================================================
\section{Ablations and analysis}
\label{sec:ablations}

Two ablations and one nuance.

\subsection{Scaffold agnosticism}

The recipe is not tied to the specific three-layer scaffold on
which the mechanism was originally characterised. An explicitly
two-layer version of the architecture (i.e. $L = 2$) trained under
the same recipe on $K_{\mathrm{hops}} = 3$ produces
target$_{N=200} = 0.992$ mean across three seeds with
QTOK-flip $\le 0.02$ on all three ($0.000$, $0.016$, $0.000$). The two-layer variant is exactly
the "degenerate" solution that $L = 3$ query-conditioned seeds fall
into under the baseline recipe \citep{fenrir2026}; when trained
explicitly, and under the filler recipe, it composes cleanly. The
recipe therefore does not depend on the three-layer scaffold as a
site of refinement; it depends only on the fast-weight write
geometry.

\subsection{Occlusion-invariance auxiliary loss}

An alternative training-time route to the same effect is a
$\lambda$-weighted KL term that penalises the change in the model's
last-position logits when a random $30\%$ of positions between the
facts and the query are replaced with SEP:
\[
  \mathcal{L}_{\mathrm{aux}} \;=\;
  \lambda \, \mathrm{KL}\!\left(
    p(\cdot \mid x)\;\|\;p(\cdot \mid \mathrm{mask}_{30\%}(x))
  \right).
\]
On $K_{\mathrm{hops}} = 3$ with $n = 1$ seed at $\lambda = 0.1$ and
$n = 2$ seeds at $\lambda = 1.0$ (plus one earlier $\lambda = 1.0$
run that failed to reach target and is excluded from the summary),
the auxiliary produces state-persistent formation with target
$\ge 0.98$ and QTOK-flip $\le 0.01$.
Replication is needed before the auxiliary sits alongside the filler
recipe as a validated intervention; we include it here as an
indication that the filler recipe is not the only training-time
route to the same effect, and that both routes share the property of
enforcing invariance to intermediate-position content.

\subsection{Filler-dependence at depth}
\label{sec:filler-dep}

The recipe's principal cost is at $K_{\mathrm{hops}} \ge 4$: mean
target$_{N=200}$ is high, but mean target$_{N=0}$ is lower, and the
gap grows with depth. On $K_{\mathrm{hops}} = 6$ passers, the mean
gap is approximately $0.43$; individual passer gaps range up to
$0.49$. Across all five filler-recipe cells, the widest single-seed
gap we observed was $0.61$ (rig 133 $s_3$ at $K_{\mathrm{hops}}=5$). The recipe biases the model toward the filler-present
distribution, and pass-seeds increasingly need filler at eval to
reach their headline accuracy. The robust property under this recipe is state-persistence and
non-divergence; balanced accuracy under both filler conditions is
not robust. Downstream use of the recipe should assume a
filler-dependence gap of roughly $0.4$ at $K_{\mathrm{hops}}=6$
(individual passers up to $0.49$), and should either evaluate
with filler or filter for seeds whose target$_{N=0}$ matches
target$_{N=200}$.

% ======================================================================
\section{Related work}
\label{sec:related}

\paragraph{Training-time interventions that shape the learned
circuit.} \citet{wang2024grokked} show that a Transformer trained on
2-hop composition transitions from a memorised solution to a
compositional circuit only via grokking, and that the training
mixture of atomic vs composed facts controls whether grokking
happens at all. Filler-augmented training is closest in spirit to
that observation: a single mixture-level intervention flips which
solution gradient descent commits to. The Transformer version acts
on in-weights composition; the fast-weight version here acts on
in-context composition. \citet{olsson2022inductionheads} identify a
similar training-time phase change for induction heads.

\paragraph{Filler / pause / think tokens.} \citet{goyal2023think}
and \citet{pfau2024let} show that inserting learnable or fixed pause
tokens between a prompt and its answer can buy the model extra
transformer forward passes for computation; \citet{hao2024coconut}
generalises this to continuous CoT vectors. Our filler tokens differ:
they are inserted at training time (not eval time) and their purpose
is not to grant the model extra computation but to force the write
geometry to be idle on filler positions. That is an unrelated
mechanism which happens to produce filler tokens as an operational
side effect.

\paragraph{Recall-capacity work on fast-weight and Mamba-family
layers.} \citet{arora2023zoology,arora2024based} establish the
capacity of single-hop associative recall on fast-weight and SSM
architectures. Our results are on the multi-hop generalisation of
that task, where the failure surface is qualitatively different: the
mechanism must \emph{form} a compositional circuit, not just
allocate memory to a set of $(k,v)$ pairs.

\paragraph{Fast-weight lineage.} The mixer that carries these
experiments is a member of the fast-weight family that runs from
\citet{schmidhuber1992learning,schmidhuber1993reducing} through
\citet{schlag2021linear}, \citet{yang2024deltanet},
\citet{yang2024gated}, \citet{peng2024rwkv7}, \citet{sun2024ttt},
\citet{behrouz2025titans}, and the self-referential branch
\citep{irie2022srwm,irie2021going}. The specific address-side
residual we use, and its multi-hop composition properties, are the
subject of a companion paper \citep{fenrir2026}. The filler recipe
reported here is expected to be architecture-generic within this
family, but we have tested it only on the specific mixer described
in Sec.~\ref{sec:background}; other members of the family may
require separate verification.

% ======================================================================
\section{Limitations}
\label{sec:limitations}

The recipe has been tested at one architecture, one task family
(injective-connectivity factorised multi-hop retrieval), one
sequence-length regime (up to $2000$ filler positions at eval), and
one hardware setup (AMD Radeon RX 7900 XTX with bf16 autocast). The strongest
empirical claims (pass rates and QTOK-flip means) come from
$n = 5$ replications per cell, which is enough to distinguish
"formation mode is steered categorically" from "formation mode has
a small drift" but not enough to characterise the tails. The
filler-dependence result reports a mean and range but not a
distribution over seeds. The two-layer ablation is $n = 3$; the
occlusion-invariance auxiliary ablation is $n = 1$ per $\lambda$
and should be treated as suggestive rather than validated. The
$K_{\mathrm{hops}} > 6$ regime is not tested; there is no
architectural reason to expect the recipe to fail at
$K_{\mathrm{hops}} = 7$ or higher, but our data cannot rule it out.

% ======================================================================
\section{Conclusion}
\label{sec:conclusion}

A single training-time change resolves three empirically distinct
failure modes of fast-weight recurrent multi-hop composition at
depth: query-conditioned formation, state divergence, and the
baseline depth wall. The change: uniform-random filler between the
last fact and the query, ramped in over a warm-up window, with a
joint loss on a filler-free probe batch and a filler-present target
batch. The recipe extends the trainable depth envelope by at least
one hop, is scaffold-agnostic, and admits a mechanistic reading in
terms of the write geometry it forces on filler tokens. Its
principal cost is a filler-dependence gap of roughly $0.4$ at
$K_{\mathrm{hops}}=6$ (individual passer gaps up to $0.5$) that
downstream users should plan around.
Further training-time interventions with the same joint property
(steering formation and suppressing divergence with one lever) are
likely to exist in the fast-weight family; the recipe reported here
is a first working example.

% ======================================================================
\bibliographystyle{plainnat}
\bibliography{references}

\end{document}
